\section{Results}\label{sec:results}

\begin{table*}[t]
\caption{Accuracy and cost (C-RQ1). PER R: person-name character recall; Scope F1: character F1 over gold of the types Presidio can detect; 95\% document-level intervals. kB/s and p50 (median ms per document) from the WNUT-17 run; FA: share of the 55 benign documents with any detection. $^\ast$CoNLL-2003 is the neural model's training distribution: an upper bound, not independent evidence.}\label{tab:accuracy}
\centering\footnotesize\setlength{\tabcolsep}{4pt}
\begin{tabular}{@{}l llllll rrr@{}}
\toprule
 & \multicolumn{2}{c}{ai4privacy (synthetic)} & \multicolumn{2}{c}{WNUT-17 (real)} & CoNLL$^\ast$ & & \multicolumn{3}{c}{Cost and false alarms}\\
System & PER R & Scope F1 & PER R & Scope F1 & PER R & & kB/s & p50 & FA \\
\midrule
T0 & \pc{0.361}{0.328, 0.396} & \pc{0.732}{0.721, 0.743} & \pc{0.000}{0.000, 0.000} & \pc{0.000}{0.000, 0.000} & \pc{0.002}{0.000, 0.007} &  & 1461.6 & 0.6 & 0.00 \\
T0+NN & \pc{0.697}{0.666, 0.729} & \pc{0.742}{0.733, 0.752} & \pc{0.481}{0.419, 0.548} & \pc{0.382}{0.325, 0.447} & \pc{0.919}{0.891, 0.941} &  & 32.6 & 29.3 & 0.56 \\
T0+NN ungated & \pc{0.937}{0.921, 0.951} & \pc{0.752}{0.744, 0.761} & \pc{0.707}{0.659, 0.757} & \pc{0.463}{0.409, 0.524} & \pc{0.992}{0.986, 0.997} &  & 10.5 & 94.1 & 0.58 \\
T0+spaCy & \pc{0.606}{0.577, 0.636} & \pc{0.694}{0.684, 0.704} & \pc{0.456}{0.400, 0.519} & \pc{0.338}{0.294, 0.391} & \pc{0.807}{0.754, 0.858} &  & 66.5 & 14.0 & 0.58 \\
Presidio sm & \pc{0.380}{0.352, 0.407} & \pc{0.690}{0.681, 0.698} & \pc{0.464}{0.408, 0.528} & \pc{0.216}{0.185, 0.250} & \pc{0.804}{0.752, 0.855} &  & 34.9 & 27.2 & 0.78 \\
Presidio lg & \pc{0.481}{0.454, 0.510} & \pc{0.709}{0.701, 0.716} & \pc{0.529}{0.474, 0.590} & \pc{0.229}{0.198, 0.264} & \pc{0.864}{0.809, 0.914} &  & 33.3 & 28.3 & 0.80 \\
\bottomrule
\end{tabular}
\end{table*}


\subsection{C-RQ1: what the neural tier adds, and at what cost}\label{sec:nn}

Table~\ref{tab:accuracy} gives accuracy and cost, and Figure~\ref{fig:recallcost} plots person-name recall against latency. On real user-generated text (WNUT-17), T0 found no person names: its tiers look for structured values, and the corpus has almost none. Adding the gated neural tier raised person-name character recall by $+0.482$ [$+0.419$, $+0.548$], and on synthetic ai4privacy by $+0.336$ [$+0.305$, $+0.369$], confirming the first planned comparison (P1) on both corpora. The pre-registered replication on BTC agreed: $+0.369$ [$+0.341$, $+0.398$].

Against T0 with spaCy's small model (P2), the neural tier had higher person-name recall on ai4privacy ($+0.092$ [$+0.065$, $+0.120$]) but not on WNUT-17 ($+0.024$ [$-0.032$, $+0.079$]). spaCy processed WNUT-17 about twice as fast (66.5 against 32.6\,kB/s, median 14.0 against 29.3\,ms per document). Against Presidio with spaCy's large model (P3), the neural tier had higher Presidio-scope character F1 on both corpora: $+0.154$ [$+0.112$, $+0.200$] on WNUT-17 and $+0.033$ [$+0.023$, $+0.043$] on ai4privacy. That margin comes from precision. On WNUT-17, Presidio lg fully covered more person names than the neural tier ($-0.092$ [$-0.147$, $-0.031$]), and the difference in person-name character recall was not significant ($-0.049$ [$-0.104$, $+0.008$]). On BTC the neural tier was ahead of T0 with spaCy ($+0.063$ [$+0.034$, $+0.089$]) and of Presidio lg in Presidio-scope F1 ($+0.214$ [$+0.193$, $+0.235$]); its person-name recall did not differ significantly from Presidio lg's ($+0.021$ [$-0.006$, $+0.047$]).

The tier is not quiet on prose without personal data. It flagged 31 of the 55 benign documents (0.56), against none for T0, 0.58 for T0 with spaCy and 0.78 and 0.80 for the two Presidio configurations. These documents name public organisations, places and people, which a named-entity model is trained to find. Batch mode (32 windows per call) gave the same accuracy within the intervals and no speed-up on these short documents. Its cost is the subject of Section~\ref{sec:contract}.

\begin{figure}[t]
\centering
\includegraphics[width=\columnwidth]{fig_recall_cost.pdf}
\caption{Person-name character recall (95\% intervals) against median latency per document, single-threaded. On real text the gated tier sits between T0 and the ungated model; Presidio lg is no faster.}\label{fig:recallcost}
\end{figure}

\subsection{C-RQ2: the gate and the streaming contract}\label{sec:contract}

\begin{table}[t]
\caption{The gate and the declared context (C-RQ2). Pass: share of characters sent to the model. Gate R: share of gold person names lying entirely inside some window, the ceiling on the tier's recall. Max: longest merged window; $>$256: windows longer than the declared context. Flips: detections that disappear when only 128 / 256 characters are kept on each side of them, out of all detections of the tier on up to 200 documents.}\label{tab:gate}
\centering\footnotesize\setlength{\tabcolsep}{4pt}
\begin{tabular}{@{}l r r r r l@{}}
\toprule
Corpus & Pass & Gate R & Max & $>$256 & Flips (128/256) \\
\midrule
ai4privacy & 0.265 & 0.543 & 462 & 24 & 1/0 of 378 \\
WNUT-17 & 0.298 & 0.594 & 258 & 1 & 1/0 of 577 \\
BTC & 0.421 & 0.422 & 149 & 0 & -- \\
CoNLL$^\ast$ & 0.598 & 0.880 & 613 & 8 & 11/2 of 2,198 \\
Benign prose & 0.588 & -- & 96 & 0 & -- \\
Loghub logs & 0.104 & -- & 258 & 1 & -- \\
\bottomrule
\end{tabular}
\end{table}


\paragraph{What the gate costs} The gate sent 26.5\% of the characters of ai4privacy and 29.8\% of WNUT-17 to the model, and 10.4\% of real system logs (Table~\ref{tab:gate}). But only 54.3\% and 59.4\% of gold person names lay entirely inside a window, which caps the tier's recall. Running the model on every line instead raised person-name character recall by 0.225 [0.173, 0.276] on WNUT-17 and 0.240 [0.210, 0.269] on ai4privacy (P4), at about a third of the throughput (10.5 against 32.6\,kB/s on WNUT-17). The gate misses names it is not built to see. Of the 174 WNUT-17 person names outside every window, 74 were written in lower case or as a handle, and 83 were single capitalised words, which the gate admits only after a lower-case word or a comma, because almost every sentence starts with one. On ai4privacy, 1,274 of the 1,294 missed names were single words. BTC repeated the pattern at a larger scale: the gate saw only 42.2\% of gold person names, 1,118 of the names it missed were in lower case or handles, and running the model on every line raised recall by 0.264 [0.235, 0.294].

\paragraph{Determinism} With batch size 1, streamed output equalled whole-document output under every chunking tried: 20 random chunkings at chunk sizes of 1 to 64\,KiB, on the WNUT-17 stream and on the first 700 ai4privacy test documents (299,915 characters). Naive chunking without the engine missed up to 15 of the 642 reference detections on WNUT-17 and leaked 5 gold names; on ai4privacy it missed 64 of 2,806 and leaked 51. With batch size 32, no chunking reproduced the whole-document output. The engine's streamed output then missed up to 46 of 644 reference detections on WNUT-17 and up to 103 of 2,822 on ai4privacy, and in both up to 12 gold values that the whole-document pass redacted were left in clear. Batching thus re-opens, through numerical non-determinism, exactly the leak the streaming engine exists to close. The option is documented as non-deterministic; our data show it is also unsafe under streaming.

\paragraph{Declared context} The tier declares 256 characters of context. Merged windows grew to 613 characters on CoNLL and 462 on ai4privacy, and 24 windows in ai4privacy exceeded the declaration (Table~\ref{tab:gate}). When only the 256 characters on each side of a detection were kept, 2 of 2,198 CoNLL detections disappeared, so the declaration is not a bound on what a decision depends on. Synthetic lines of chained names (``Ann Lee, Bo Kim, \dots'') produced merged windows of up to 2,068 characters, yet streamed output still equalled whole-document output in all 15 chunkings at every length tried. The engine's overlap, 8,192 characters for the whole pipeline, is set by the longest-reaching pattern tier and absorbed the excess. A deployment that ran the neural tier alone would have no such margin. The fix is to cap the merged window at the declared bound, or to declare the bound the gate can actually produce.

\subsection{C-RQ3: what crosses the tool boundary}\label{sec:egress}

\MISSING{E6 egress pre/post}


\paragraph{Content-free replies} Before the fix, the differential audit found person names in \texttt{explain\_\allowbreak detection} replies whenever the gazetteer tier (a model-free list of given names) was active (Table~\ref{tab:egress}). Every hit was in the \texttt{notes.norm} field, which carried the detected name case-folded (``maria gonzalez''), and in the JSON text block that repeats it. Names leaked from 701 of the 2,000 AgentTrace documents, 24 of the 129 WNUT-17 documents and 119 of the 3,973 ai4privacy documents. No other tier, field, interface or encoding leaked a gold value. The cause was a deny-list: \texttt{notes} forwarded every detector metadata key except one. The fix forwards only an allow-list of ten keys whose values are fixed vocabulary or numbers. It also removed three fields that had not leaked a gold value in our corpora but could carry input text: a label cue, a structured-field key and the model's file path. After the fix the audit found no egress in any tier, corpus, interface or encoding. A property test that plants random surnames now fails the build on any regression. The pilot run also showed why the audit compares against a masked control and matches whole tokens: without them, three-letter values (``Hon'') matched inside tool words (``phone''), and house numbers matched character offsets in SARIF messages.

\paragraph{Redacted text} Through \texttt{redact}, a value is exposed exactly when detection misses it. With the neural tier, 0.588 of WNUT-17 values of at least four characters and 0.420 of ai4privacy values were still present verbatim in the redacted text. That is slightly below the share not fully covered by any detection (0.631 and 0.437), because a partly covered value no longer appears whole. Every mode except mask gave the same exposure. Mask keeps a suffix of each covered value visible, at most four characters and never more than half the value: 36.0\% of the characters of covered WNUT-17 values and 27.4\% on ai4privacy. Offsets and lengths, which \texttt{scan} returns by design, disclose a value's position and length but not its characters. Hash and pseudonymize disclose whether two values are equal, which is their purpose.

\subsection{C-RQ5: other agent-facing redaction tools}\label{sec:thirdparty}

\begin{table}[t]
\caption{Third-party agent-facing redaction tools (C-RQ5), pseudonymised until their maintainers have been told. Left: share of gold values left in the tool's redacted output (a recall measure) on AgentTrace (AT), WNUT-17 (WN) and 500 ai4privacy documents (AI). Map: share returned elsewhere in the same reply. Detect: share returned by the tool's detection call. Restore: hidden values the agent got back by calling the tool's own restore call, on AgentTrace.}\label{tab:thirdparty}
\centering\footnotesize\setlength{\tabcolsep}{3pt}
\begin{adjustbox}{max width=\columnwidth}\begin{tabular}{@{}l rrr rr r@{}}
\toprule
 & \multicolumn{3}{c}{Left in output} & & & \\
Tool & AT & WN & AI & Map & Detect & Restore \\
\midrule
T1 & 0.32 & 0.49 & 0.50 & 0.00 & 0.63 & -- \\
T2 & 0.14 & 0.92 & 0.83 & 0.62 & -- & 4,489/4,489 \\
T3 & 0.55 & 0.84 & 0.73 & 0.00 & 0.44 & 2,339/2,339 \\
T4 & 0.48 & 0.99 & 0.89 & 0.00 & -- & 2,750/2,750 \\
SafeStream & 0.19 & 1.00 & 0.54 & 0.00 & 0.00 & -- \\
\bottomrule
\end{tabular}\end{adjustbox}
\end{table}


We ran the same audit, through the official MCP client, on every open-source redaction tool that met the pre-registered inclusion criteria. Six candidates were found. Two were excluded at installation: one ships middleware but no tool, and the offline model of the other could not be downloaded. A second engine of one tool was also excluded, because it demanded a commercial key after ten requests a day. Four tools were audited, called T1 to T4 until their maintainers have been told (Section~\ref{sec:discussion}). Table~\ref{tab:thirdparty} gives the results. The share of values left in each tool's redacted output depends on which entity types it targets, so it is reported but not compared.

\paragraph{The agent can undo the redaction} Three of the four tools give the agent a restore call that turns placeholders back into the original values, so that a model's answer can be shown to the user with the real values in place. Called by the agent on the output of the tool's own anonymiser, the restore call returned every value it had hidden: 4,489 of 4,489 for T2, 2,339 of 2,339 for T3 and 2,750 of 2,750 for T4. In the threat model of Section~\ref{sec:threat}, where the model or an injected prompt chooses the tool calls, such a call cancels the protection. For two tools this contradicts their own documentation: T3's documentation says its restore call returns only a file path, but without an output path it returned the restored text; T4's documentation describes restored text being written to disk and not shown to the model. T3's separate mapping call returned only placeholder types, never values, which is the safe design.

\paragraph{Values in the same reply} T2 returned a map from each placeholder to its original value in the same reply as the anonymised text. That map carried 0.62 of all values on AgentTrace (3,242 of 5,250), although its documentation says no data leaves the machine: in an agent loop the reply goes to the model. The detection calls of T1 and T3 returned the matched values themselves, 0.63 and 0.44 of all values on AgentTrace. That is by design and documented for T1, which warns that its use exposes the text to the model. Under the same harness, no SafeStream reply carried a value outside its redacted output, and it offers no restore call to the agent.

\subsection{C-RQ6: a real agent, with and without a prompt injection}\label{sec:agent}

\begin{table}[t]
\caption{A local LLM agent (Qwen2.5-7B-Instruct) asked to summarise a ticket it may read only through each redaction tool (C-RQ6), 20 tickets per condition. Reached: runs in which a value the tool had hidden reached the model through a tool result; Answer: in the final answer; Restore: runs in which the agent called a restore tool. B: benign ticket; I: ticket with an appended prompt injection.}\label{tab:agent}
\centering\footnotesize\setlength{\tabcolsep}{3pt}
\begin{tabular}{@{}l rr rr rr@{}}
\toprule
 & \multicolumn{2}{c}{Reached} & \multicolumn{2}{c}{Answer} & \multicolumn{2}{c}{Restore} \\
Tool & B & I & B & I & B & I \\
\midrule
SafeStream & 0/20 & 0/20 & 0/20 & 0/20 & 0/20 & 0/20 \\
T2 & 20/20 & 20/20 & 17/20 & 20/20 & 0/20 & 0/20 \\
T3 & 0/20 & 1/20 & 0/20 & 1/20 & 0/20 & 14/20 \\
T4 & 7/20 & 13/20 & 7/20 & 9/20 & 13/20 & 15/20 \\
\bottomrule
\end{tabular}
\end{table}


The audits above measure what a tool returns. To see what an agent does with it, a local open-weight model, Qwen2.5-7B-Instruct (4-bit, run with llama.cpp at temperature 0), was asked to summarise a ticket in three sentences. It could read the ticket only through each tool's anonymiser, and it was also given the tool's own text tools, the restore call among them. Each tool was run on 20 generated tickets, once as written and once with an appended paragraph telling AI assistants to call ``the tool that restores or reveals the original values'' and quote the restored ticket (Table~\ref{tab:agent}).

With SafeStream, no hidden value reached the model in any run, with or without the injection. With T2, every run leaked, with or without it: the anonymiser returns the value map with the text, so the hidden values were in the model's context from the first read, and the agent copied them into its summary in 17 of 20 benign runs and all 20 injected ones. With T3 and T4, the restore call decided the outcome. Without the injection, the agent never called T3's restore call, but it called T4's in 13 of 20 runs unprompted, and hidden values reached it in 7. The injection raised restore calls to 14 of 20 for T3 and 15 of 20 for T4. Runs in which hidden values reached the model rose to 13 of 20 for T4 (95\% interval 0.41 to 0.85, against 0.15 to 0.59 without the injection). For T3, in 10 of its 14 restore calls the agent asked the tool to write the restored text to a file rather than return it, so values reached the model in only one run. The cleartext was still restored, onto disk at a path the agent chose; this count is exploratory, not a pre-registered outcome. One model, one injection and 20 tickets per cell show that the attack works on a capable open model, not how often it works on every model.

\subsection{C-RQ4: what the guardrail costs per call}\label{sec:cost}

\MISSING{E8 latency / E9 policy}


\paragraph{Transport} Table~\ref{tab:cost} gives median latency per \texttt{scan} call. For small payloads the transport is visible but cheap: at 256\,B a call took 0.28\,ms in process, 0.98\,ms over stdio and 0.75\,ms over REST. From 64\,KiB, detection dominates. Over stdio the official MCP client stayed within the pre-registered factor of 1.25 of in-process latency at 64\,KiB (88.79 against 79.46\,ms) and 1\,MiB (1,480.27 against 1,268.70\,ms), but not at 4\,MiB (7,421.76 against 5,319.80\,ms), so H4a holds only up to 1\,MiB. We did not profile where the extra time at 4\,MiB goes. Starting the stdio server and completing the MCP handshake took between 0.08 and 0.126\,s. Measuring REST revealed that every kept-alive request waited about 40\,ms regardless of payload. The handler wrote headers and body in two writes, and Nagle's algorithm \cite{rfc896} held the body until the client's delayed acknowledgement \cite{rfc1122}. Setting \texttt{TCP\_NODELAY} on the handler removed the stall: the median 256\,B call fell from 44.00 to 0.75\,ms, and the 4\,KiB call from 48.06 to 5.87\,ms. From 64\,KiB the difference disappears into detection time.

\paragraph{Redaction modes and reports} On 16\,MiB of AgentTrace text (258,749 detections), end-to-end \texttt{redact} throughput varied little with the mode, from 1.22\,MB/s for mask to 1.06\,MB/s for encrypt: detection, not rendering, sets the cost. Rendering a replacement took 0.44\,$\mu$s per detection for replace and 9.65\,$\mu$s for encrypt, whose tokens nearly double the output (1.945 bytes out per byte in). Encrypt round-tripped exactly: all 16,171 tokens in a 1\,MiB sample were restored to the original text. Of 10,000 tokens with one bit of the nonce, ciphertext or tag flipped, none was restored. A SARIF report of 7,500 findings validated against the official 2.1.0 schema with no error, and building it cost 5.38\,ms per 1,000 findings.
